\documentclass[11pt]{article}
\usepackage[margin=1in]{geometry}
\usepackage{amsmath,amssymb,booktabs,microtype,graphicx,xcolor}
\usepackage[round]{natbib}
\usepackage[colorlinks=true,linkcolor=blue!50!black,urlcolor=blue!50!black,citecolor=blue!50!black]{hyperref}
\setlength{\parskip}{0.35em}\setlength{\parindent}{0em}
\title{\textbf{Margin Is Not a Footnote:\\ Reg-T Feasibility as a First-Class Backtest Output}}
\author{Charlie Yan\thanks{Draft v0.3 (journal-length), August 2026. Margin figures come from
our own Reg-T implementation, computed daily from the tape's actual positions (strikes, widths,
credits, spot), hand-verified on sampled days against the published formulas, with two
independent implementations agreeing to $\sim$1\%. They are \emph{our arithmetic, not the
broker's}: validation against broker-computed requirements is pending and is flagged throughout.
Requirements are reported as percentages of a reference net liquidation value; because
requirements scale with position count, these ratios are scale-invariant and the normalization
is a reference figure rather than deployed capital. Performance basis where quoted: full cross
fills, ex-commission, marked to market daily on the full calendar, 2020-09 to 2026-05.
Code, pre-registration documents and a data-availability statement: \url{https://github.com/charlieyanhx/margin-first-class}.}}
\date{August 2026}
\begin{document}\maketitle

\begin{abstract}
A backtest of a margined strategy that reports a Sharpe ratio without a margin path is
incomplete in a way that can invert its conclusion. We document a short-premium index options
book whose uncapped configuration backtests attractively on every performance statistic---and
whose daily Reg-T requirement, reconstructed from the tape's own positions, \emph{peaks at
151--154\% of the account's net liquidation value}. The strategy is infeasible: it breaches its
own capital on precisely the stress day that generates its risk premium. The naked short-call
sleeve alone peaks at 108--111\% of the whole account. Adding a far wing to that sleeve---an
action costing roughly one-fifteenth of the equivalent put wing, because of smile
asymmetry---halves the book's peak to 77\% and converts an infeasible strategy into a
deployable one at approximately zero cost to expected P\&L. None of this is visible in any
performance statistic, because margin peaks are joint events: entries cluster, holds stack, and
per-position requirements widen simultaneously on the same stress day, so the peak-to-median
ratio runs 6.7 to 16.1$\times$ across configurations while every performance statistic is a
moment of a single daily series. We propose a minimum reporting standard for margined
backtests---the dated requirement path, the peak's anatomy, an intrinsic-bound gap table on
actual daily inventory, and return on deployed margin alongside return on capital---and we
document three portable findings encountered in building it: a stale measurement window
understated the true peak by half; iron-condor netting cuts the peak requirement 15\% with
identical positions and P\&L, bounded by expiry alignment; and at small scale collateral is
usually not the binding constraint (signal supply is), so margin optimization pays exactly at
the scale-up where it stops being optional.
\end{abstract}

\section{The statistic nobody reports}
Consider two strategies with identical Sharpe ratios, identical maximum drawdowns, identical
turnover, and identical per-year P\&L. One of them can be run; the other cannot, because on a
single day in April 2025 its collateral requirement exceeds the account that is running it. No
statistic in the standard performance toolkit distinguishes them. That gap is this paper's
subject.

The omission is not a matter of degree. For an unlevered long-only strategy, capital
consumption is trivially the position value, and reporting it adds nothing. For a margined
short-option book it is a separate stochastic process, driven by different variables than P\&L,
peaking at different times, and capable of terminating the strategy independently of its
returns. A backtest that omits it has not modelled the strategy; it has modelled the strategy's
profits conditional on the strategy being permitted to continue---a condition it never checked.

We make three contributions. Section~\ref{sec:case} presents a worked case in which the margin
path inverts the deployment conclusion for a book that every performance statistic approves.
Section~\ref{sec:why} explains why performance statistics are structurally blind to this,
through the joint-event mechanism, and quantifies the blindness. Section~\ref{sec:standard}
proposes a reporting standard that costs a few lines of code against a position file, plus
three findings from building it that we expect to generalize.

\section{The case}\label{sec:case}
The book comprises three sleeves of index option spreads---short put verticals, short call
positions, and a delta-hedged relative-value sleeve---over 2020--2026, with entries clustering
in volatility episodes and multi-day holding periods. Its backtested performance was acceptable
on every statistic the standard toolkit produces, computed on honest fills and daily marking.

Reconstructing the daily Reg-T requirement from the tape's actual positions, with two
independent implementations agreeing to $\sim$1\% and anchored on hand-verified days:

\begin{table}[h]\centering\small
\begin{tabular}{lcccc}\toprule
Configuration (\% of reference NLV) & Median & Mean & P90 & \textbf{Peak (date)} \\ \midrule
Put sleeve, winged & 5.5 & 7.9 & 17.5 & 44.1\% \;(2025-04-15) \\
Call sleeve, \textbf{naked (as designed)} & 11.2 & 15.5 & 38 & \textbf{108--111\%} \;(2025-04-09) \\
Call sleeve, winged & 2.2 & 3.5 & 7.8 & 35.4\% \;(2025-04-08) \\
\midrule
\textbf{Book, naked-call configuration} & 22 & --- & --- & \textbf{151--154\%} \;(2025-04-09) \\
\textbf{Book, winged configuration} & 11.5 & --- & --- & \textbf{76.9\%} \\
\bottomrule\end{tabular}
\caption{Audited daily Reg-T requirement by configuration. The naked-call book breaches its own
account; the winged book does not.}
\label{tab:margin}\end{table}
\begin{figure}[h]\centering
\includegraphics[width=\textwidth]{fig_p5_margin.pdf}
\caption{Exhibit 1. Left: the requirement distribution by configuration, against the account's
capital line. Right: the peak-to-median ratio---what a ``typical utilization'' figure conceals.}
\label{fig:margin}\end{figure}

The breach day is not exotic. It is the April 2025 stress episode: roughly one hundred open
short calls near the money after a $+9.4\%$ melt-up, at a hand-verified requirement per
contract in the five figures. The strategy's premium is compensation for precisely this kind of
day; the margin system's response to precisely this kind of day is what makes the strategy
impossible to hold through it. That is not a coincidence but the structure of the problem: the
event that justifies the premium and the event that triggers the collateral call are the same
event.

One presentational recommendation follows directly. Plotting the margin path \emph{against the
equity curve}, rather than alone, makes legible the relationship between what a strategy has
earned to date and what it must post today---a relationship that is unfavourable early in any
strategy's life, since requirements scale with position count from day one while cumulative
P\&L starts at zero. We know of no performance statistic that substitutes for looking at the
two series together.

\section{Why performance statistics cannot see this}\label{sec:why}
Margin peaks are \emph{joint} events. On the same stress day three things happen at once:
entries have clustered, because the signal fires repeatedly in episodes; holds have stacked,
because positions from the preceding weeks are all still open; and each position's requirement
has widened, because strikes have moved toward the money and per-position spread widths roughly
doubled through the episode. The margin path is therefore the product of three correlated
stress responses, while every performance statistic---mean, volatility, Sharpe ratio, even
maximum drawdown---is a moment or an extremum of a \emph{single} daily P\&L series. The
information is not compressed by those statistics; it is absent from their inputs.

The quantitative consequence is the peak-to-median ratio, which runs 6.7 to 16.1$\times$ across
the configurations in Table~\ref{tab:margin} (Exhibit~1, right panel). Typical utilization is
11.5\% of capital for the winged configuration and 22\% for the naked one (16\% median in the
small-account variant): figures that, reported alone, are not merely incomplete but actively
misleading, since each is roughly an order of magnitude below the number that determines
whether the strategy can be run. Reporting the median without the peak is the collateral analogue of
reporting a mean return without a drawdown.

Three portable specifics emerged from the audit.

\textbf{(a) The stale-window trap.} A previously recorded figure of ``129\% peak at production
sizing'' turned out to be the peak of a truncated early window---election week 2020---not
exceeded again until 2022, after which position concurrency tripled. Measured on the full
sample, the true peak was roughly double. Margin paths must span the full sample for exactly
the reason performance must, and for an additional one: concurrency is a slow-moving structural
variable, so an early-window margin measurement is not merely noisy but systematically biased
low for any strategy whose position count grew.

\textbf{(b) The cheap wing.} Because of put-call smile asymmetry, the far call wing that
converts the naked sleeve to defined risk costs roughly one-fifteenth of the equivalent put
wing, while halving the book's peak requirement---the cheapest unit of feasibility in the book,
at approximately zero cost to expected P\&L. This reframes what wings are for. In this book the
put wing recovers only 22\% of worst-case blowup losses, so as insurance it is poor value; as
collateral relief it is decisive. \emph{Wings are structure, not insurance}, and their margin
function dominates their hedging function in the deployment decision.

\textbf{(c) Netting.} Pairing put and call spreads of matching expiry into iron condors---
identical positions, identical P\&L, since Reg-T charges the greater side rather than the
sum---cuts the peak requirement 15\% and the median 13\%. The binding constraint is expiry
alignment: the two sleeves choose expiries independently and coincide only 31\% of the time.
Deliberately aligning them should push the netting rate toward its limit at near-zero strategy
cost, a one-variable change we have not yet tested. Broker recognition of the pairing must be
confirmed before any of this is relied upon.

\section{The gap table, and why history understates it}\label{sec:gap}
Realized drawdown history is the wrong basis for sizing a short-premium book whenever the
sample lacks a generational stress day. Ours spans 2020-09 onward and contains neither March
2020 nor February 2018---a fact easy to state and easy to forget when a five-year backtest
looks long. We therefore report, alongside the margin path, an \emph{overnight gap table}: for
every day's actual inventory, the position-level intrinsic-value bound on the overnight loss
for gaps out to $\pm 60\%$.

The construction is deliberately model-free. It asks what each open position is worth at
expiry-equivalent intrinsic value under a hypothesised gap, sums across the actual inventory of
that day, and takes the worst day. It requires no volatility model, no correlation assumption,
and no distributional claim---which is what makes it a bound rather than a forecast. On the
studied small-account configuration a $-15\%$ overnight gap bounds the loss at 49\% of the
account on the worst inventory day, and no gap on any day in the sample reaches forced
liquidation: a statement made possible only by the defined-risk wings, and one that no
realized-history statistic can make. Against that bound, the realized maximum drawdown
understates the structural tail by roughly six times.

We regard the gap table as the natural companion to the margin path because the two answer the
same question from opposite directions. The margin path asks: on the worst day the sample
actually contained, could the strategy post collateral? The gap table asks: on a worse day the
sample did not contain, would the strategy survive at all?

\section{Return on margin versus return on capital}\label{sec:ronm}
The book's return on margin actually deployed is roughly 111\% per year; its return on total
capital is an order of magnitude lower. The gap is not inefficiency. The idle capital is the
reserve that makes the peak days survivable, and its opportunity cost is properly a cost of the
strategy, not a shortfall in its execution.

Reporting both numbers prices the reserve explicitly. Reporting only return on capital hides
the strategy's true capital intensity and makes it look worse than the trading is; reporting
only return on deployed margin overstates scalability by pretending the reserve is free. The
pair, together, is what a capital allocator needs, and the difference between them is the
quantity a portfolio-margin regime or a netting improvement actually attacks.

One caution against over-optimizing here. At small scale, collateral is usually \emph{not} the
binding constraint in this book---median utilization is 16\%, and what binds is signal supply,
roughly two entries per day. Margin optimization therefore pays at the scale-up where the peak
becomes binding, which is exactly the point at which it stops being optional. Work on it before
then and it is a solution without a problem; discover it after and the strategy has already
been sized on a false premise.

\section{A proposed reporting standard}\label{sec:standard}
For any margined backtest we suggest reporting, as first-class outputs alongside performance:
\begin{enumerate}\itemsep0.05em
\item \textbf{The requirement path}: median, 90th percentile and peak, each dated, over the full
sample.
\item \textbf{The peak's anatomy}: concurrency, per-position widening, and entry clustering on
the peak day---the three multipliers, separated.
\item \textbf{A feasibility statement}: peak as a percentage of capital, with the margin regime
named (Reg-T or portfolio margin) and the conclusion's sensitivity to that choice stated.
\item \textbf{An intrinsic-bound gap table} on actual daily inventory, not a distributional
tail estimate.
\item \textbf{Return on deployed margin next to return on capital}, with the reserve priced.
\item \textbf{Where infeasible, the cheapest structural remedy} and its P\&L cost---in our case,
a wing costing approximately nothing that halved the peak.
\end{enumerate}
Items 1--3 require a few lines of code against the position file a backtest already produces;
none needs broker connectivity, though broker validation is the final word and ours is pending.
We would go further and suggest that a margined backtest reporting a Sharpe ratio without at
least item~1 should be regarded the way a backtest without transaction costs is now regarded:
not wrong, exactly, but not yet a result.

\section{Relation to prior work}
Margin and leverage constraints are central to modern asset pricing. \citet{fp} build the
betting-against-beta factor on the premise that leverage constraints bind and are priced;
\citet{bp} model the funding-liquidity spiral by which margin requirements and asset prices
reinforce each other in stress---the macro version of the mechanism that makes our peak
coincide with our premium event. \citet{gk} formalize collateral-constrained equilibrium; the
margin-based asset-pricing literature that follows treats the requirement as a state variable
rather than an afterthought.

Closest to our subject, \citet{hitz} show that equity option returns carry a \emph{margin
premium} determined by an option's contribution to dealers' overall requirements---margin as a
priced state variable. That is the mirror of our question: they ask what margin does to
expected returns, we ask what it does to feasibility. \citet{ss} is nearer still in spirit,
documenting that attractive option strategies can be undone by margin calls.

\textbf{What is and is not novel here.} We initially believed margin reporting to be absent
from option-strategy backtesting; that overstates the case and we correct it. Retail and
vendor backtest publishers routinely report capital utilization, and some size starting capital
\emph{to} the margin peak (targeting 80--100\% maximum utilization), which is a stronger
practice than merely reporting a median; return-on-margin alongside notional return is likewise
established in that literature. Institutional risk systems compute requirements continuously.
What we have not found, and what we claim, is narrower and we state it precisely: (i) a worked
case in which the margin path \emph{inverts} a deployment conclusion that every performance
statistic approves; (ii) the joint-event decomposition explaining why performance statistics are
structurally blind to it, quantified as a 6.7--16.1$\times$ peak-to-median ratio; (iii) the
intrinsic-bound gap table on actual daily inventory as a complement to realized-drawdown
history; and (iv) the specific empirical result that smile asymmetry makes the feasibility-
restoring wing cost roughly one-fifteenth of its put-side equivalent. Readers who already size
to peak utilization have item (i) covered; items (ii)--(iv) we believe are new.

\section{Limitations}
Our Reg-T implementation is hand-verified against published formulas and cross-checked between
two independent builds, but it is not yet validated against broker-computed requirements;
house minimums, add-ons and portfolio-margin models can all bind earlier than our arithmetic
suggests, and the direction of that error is unfavourable. The portfolio-margin figures in the
underlying audit are intrinsic-only approximations. The case is one book, one underlying, one
margin regime, and a sample lacking a generational stress event---the last being precisely why
Section~\ref{sec:gap} exists. The proposed standard, however, does not depend on the case's
particulars: any strategy whose collateral requirement is a separate stochastic process from
its P\&L needs the requirement reported, and most margined strategies qualify.

\begin{thebibliography}{99}\small
\bibitem[Brunnermeier and Pedersen(2009)]{bp} Brunnermeier, M., Pedersen, L. (2009). Market
liquidity and funding liquidity. \emph{Review of Financial Studies} 22(6).
\bibitem[Frazzini and Pedersen(2014)]{fp} Frazzini, A., Pedersen, L. (2014). Betting against
beta. \emph{Journal of Financial Economics} 111(1).
\bibitem[Geanakoplos(2010)]{gk} Geanakoplos, J. (2010). The leverage cycle. \emph{NBER
Macroeconomics Annual} 24(1).
\bibitem[Hitzemann et al.(2021)]{hitz} Hitzemann, S., Hofmann, M., Uhrig-Homburg, M., Wagner,
C. (2021). Margin requirements and equity option returns. Working paper, SSRN 2789113.
\bibitem[Santa-Clara and Saretto(2009)]{ss} Santa-Clara, P., Saretto, A. (2009). Option
strategies: good deals and margin calls. \emph{Journal of Financial Markets} 12(3).

\end{thebibliography}
\end{document}
